% command to make sure colorboxes don't create whitespace between code lines
\fboxsep=.6pt

% default settings for listings
\lstset{style=python, basicstyle=\scriptsize\ttfamily,escapechar=!}

\begin{frame}
    \centering
    \begin{figure}[H]
        \includegraphics[keepaspectratio, width=0.7\paperwidth, height=.6\paperheight]{Logos/FreeFlower.pdf}
    \end{figure}
    Informatik: Programmieren\\
    \emoji{abacus} Kapitel 03: \textbf{Werte bei der Ausführung mitgeben}
\end{frame}



\begin{frame}[fragile]{Vieleck mit Variablen als \textit{direkt gesetzt} \emoji{handshake}}
\begin{columns}
\begin{column}{.7\textwidth}
\begin{lstlisting}
import turtle as t

# Anzahl Ecken und Seitenlänge
ecken = 5
laenge = 100

for _ in range(ecken):
    t.forward(laenge)
    t.right(360/ecken)
\end{lstlisting}
$\rightarrow$ Einfach, direkt!
\end{column}
\begin{column}{.3\textwidth}
\begin{figure}
\centering
\includegraphics[width=\linewidth]{Grundlagen_Info/00_Programmieren/Figures/vieleck}
\end{figure}
\end{column}
\end{columns}
\end{frame}

\begin{frame}[fragile]{Vieleck mit \textit{hard-coded} Variablen \emoji{pencil}}
\begin{columns}
\begin{column}{.7\textwidth}
\begin{lstlisting}
import turtle as t

# Anzahl Ecken und Seitenlänge
# sind direkt im Code
ecken = 5
laenge = 100

for _ in range(ecken):
    t.forward(laenge)
    t.right(360/ecken)
\end{lstlisting}
$\rightarrow$ Variablen sind fest!
\end{column}
\begin{column}{.3\textwidth}
\begin{figure}
\centering
\includegraphics[width=\linewidth]{Grundlagen_Info/00_Programmieren/Figures/vieleck}
\end{figure}
\end{column}
\end{columns}
\end{frame}


\begin{frame}[fragile]{Vieleck mit Variablen, \textit{on-the-fly} \emoji{airplane}}{Was passiert?}
\begin{lstlisting}
import turtle as t

# Anzahl Ecken und Seitenlänge
# mit input('...')
ecken = input("Wie viele Ecken?")
laenge = input("Seitenlänge?")

for _ in range(ecken):
    t.forward(laenge)
    t.right(360/ecken)
\end{lstlisting}
\pause
\textcolor{red}{\texttt{TypeError: 'str' object cannot be interpreted as an integer}}
\begin{itemize}
\item \lstinline|input| gibt immer einen Text (\lstinline|str|) zurück, z.B. \lstinline|"5"| statt \lstinline|5|
\item Mit einem Text kann Python nicht rechnen und keine Schleife wiederholen
\end{itemize}
\end{frame}

\begin{frame}[fragile]{Vieleck mit Variablen, \textit{on-the-fly} \emoji{airplane}}{Text in Zahl umwandeln mit \lstinline|int(...)|}
\begin{columns}
\begin{column}{.7\textwidth}
\begin{lstlisting}
import turtle as t

# Anzahl Ecken und Seitenlänge
# mit input('...') als ganze Zahl
ecken = int(input("Wie viele Ecken?"))
laenge = int(input("Seitenlänge?"))

for _ in range(ecken):
    t.forward(laenge)
    t.right(360/ecken)
\end{lstlisting}
\end{column}
\begin{column}{.3\textwidth}
\begin{figure}
\centering
\includegraphics[width=\linewidth]{Grundlagen_Info/00_Programmieren/Figures/vieleck}
\end{figure}
\end{column}
\end{columns}
\end{frame}

\begin{frame}[fragile]{Vieleck mit Variablen, \textit{on-the-fly} \emoji{airplane}}{Mehrere Vielecke}
\begin{lstlisting}
import turtle as t

# Drei Vielecke wiederholt mit input
for _ in range(3):
    ecken = int(input("Wie viele Ecken?"))
    laenge = int(input("Seitenlänge?"))
    for _ in range(ecken):
        t.forward(laenge)
        t.right(360/ecken)
\end{lstlisting}
$\rightarrow$ Demo
\end{frame}

\begin{frame}[fragile]{Zusammenfassung}
\begin{itemize}
\item Bei \lstinline|input| sind die Variablen-Werte vor der Programm-Ausführung \textit{unbekannt}!
\item Ermöglicht es, flexibel Variablenwerte \textit{während} der Programm-Ausführung zu definieren (\textit{on-the-fly} \emoji{airplane})
\end{itemize}
\end{frame}

\begin{frame}[fragile]{Auftrag}
\begin{itemize}
\item \abgabebox{Aufgabe 3.17}
\end{itemize}
\end{frame}

% Kurzinput zu Skript Abschnitt «Debugging», damit die Aufgaben 3.18-3.19 hier folgen können
\begin{frame}[fragile]{Debugging: Arten von Fehlern}
\begin{itemize}
\item \textbf{Syntaxfehler}: Schreibweise nicht erlaubt, das Programm startet nicht.\\
\lstinline|print "Hello World"|
\item \textbf{Laufzeitfehler}: Fehler erst bei der Ausführung, z.\,B. ein Name, den es nicht gibt.\\
\lstinline|X = 10 + 5| und danach \lstinline|print(x)|
\item \textbf{Semantische Fehler}: keine Fehlermeldung, aber ein falsches Resultat.\\
\lstinline|mittelwert = 3 + 13 / 2| ergibt 9.5 statt 8
\end{itemize}
\textbf{Strategien}: Fehlermeldung und Zeilennummer lesen, Zwischenresultate mit \lstinline|print| ausgeben, Beispiele von Hand nachrechnen, mit besonderen Werten testen.
\end{frame}

\begin{frame}{Auftrag}
\begin{itemize}
\item \aufgabebox{Aufgaben 3.18-3.19}
\item \challengebox{Aufgaben 3.20-3.22}
\end{itemize}
\end{frame}